\documentclass[12pt,a4paper]{article}

\usepackage[margin=1in]{geometry}
\usepackage{fontspec}
\usepackage{xcolor}
\usepackage{hyperref}
\usepackage{enumitem}
\usepackage{booktabs}
\usepackage{listings}
\usepackage{float}

\defaultfontfeatures{Ligatures=TeX,Renderer=HarfBuzz}
\directlua{
  luaotfload.add_fallback("banglafallback", {
    "Noto Serif Bengali:mode=harf;script=bng2;language=dflt"
  })
  luaotfload.add_fallback("banglamonofallback", {
    "Noto Sans Bengali:mode=harf;script=bng2;language=dflt"
  })
}
\setmainfont[RawFeature={fallback=banglafallback}]{Latin Modern Roman}
\setmonofont[RawFeature={fallback=banglamonofallback}]{Latin Modern Mono}
\newfontfamily\banglafont[
  RawFeature={fallback=banglafallback}
]{Latin Modern Roman}
\newfontfamily\banglasans[
  Script=Bengali,
  Renderer=HarfBuzz
]{Noto Sans Bengali}
\newfontfamily\banglamono[
  Script=Bengali,
  Renderer=HarfBuzz
]{Noto Sans Bengali}

\newcommand{\bn}[1]{{\normalfont\banglafont #1}}
\newcommand{\bnsf}[1]{{\normalfont\banglasans #1}}
\newcommand{\code}[1]{\texttt{#1}}

\pdfstringdefDisableCommands{%
  \def\bn#1{#1}%
  \def\bnsf#1{#1}%
}

\hypersetup{
  colorlinks=true,
  linkcolor=black,
  urlcolor=blue
}

\title{\bn{LaTeX Basic Command Guide}}
\author{\bn{বাংলা ব্যাখ্যাসহ}}
\date{\today}

\begin{document}

\maketitle
\tableofcontents
\newpage

% ---------------------------------------------
% SECTION 1 (NEW): Document Structure
% ---------------------------------------------
\section{\bn{Document Structure (Preamble)}}

\bn{একটি LaTeX file সাধারণত দুই ভাগে ভাগ করা হয়:}

\begin{enumerate}
  \item \bn{\textbf{Preamble} --- `\textbackslash documentclass` থেকে `\textbackslash begin\{document\}` পর্যন্ত। এখানে package load করা হয় এবং global setting দেওয়া হয়।}
  \item \bn{\textbf{Body} --- `\textbackslash begin\{document\}` থেকে `\textbackslash end\{document\}` পর্যন্ত। এখানে মূল লেখা থাকে।}
\end{enumerate}

\begin{verbatim}
\documentclass[12pt,a4paper]{article}  % document type

% -- Packages --
\usepackage[margin=1in]{geometry}      % page margin
\usepackage{graphicx}                  % image support
\usepackage{amsmath}                   % math support
\usepackage{xcolor}                    % color support
\usepackage{hyperref}                  % hyperlink support
\usepackage{booktabs}                  % nice table lines
\usepackage{float}                     % figure placement [H]

% -- Title Info --
\title{My Document}
\author{Your Name}
\date{\today}

\begin{document}

\maketitle           % title print করে
\tableofcontents     % table of contents তৈরি করে
\newpage             % নতুন পেজে যায়

% এখানে মূল লেখা লিখুন ...

\end{document}
\end{verbatim}

\bn{`\textbackslash documentclass` option-এ `12pt` মানে font size 12, `a4paper` মানে A4 size page।}

% ---------------------------------------------
% SECTION 2: How LaTeX works
% ---------------------------------------------
\section{\bn{LaTeX কীভাবে কাজ করে}}

\bn{LaTeX হলো একটি document preparation system. এখানে Microsoft Word-এর মতো
সরাসরি ডিজাইন না করে আপনি source code লিখেন। তারপর LaTeX সেই code compile করে
একটি সুন্দর PDF বানায়।}

\begin{center}
\bn{.tex file + image + bibliography} $\rightarrow$ \bn{compile} $\rightarrow$ \bn{PDF}
\end{center}

\bn{সাধারণভাবে কাজ করার ধাপ:}

\begin{enumerate}
  \item \bn{একটি `.tex` file-এ লেখা লিখুন।}
  \item \bn{প্রয়োজন হলে image `figures/` folder-এ রাখুন।}
  \item \bn{reference থাকলে `references.bib` file-এ দিন।}
  \item \bn{compile করুন।}
  \item \bn{শেষে PDF দেখুন।}
\end{enumerate}

% ---------------------------------------------
% SECTION 3: Command Syntax
% ---------------------------------------------
\section{\bn{Command Syntax}}

\bn{LaTeX command সাধারণত backslash দিয়ে শুরু হয়।}

\begin{verbatim}
\command{required value}
\command[optional value]{required value}
\end{verbatim}

\bn{উদাহরণ:}

\begin{verbatim}
\textbf{This text is bold}
\includegraphics[width=0.8\textwidth]{figures/image.png}
\end{verbatim}

\bn{এখানে curly braces `{}`-এর ভিতরের অংশ mandatory বা required value।
Square brackets `[]`-এর ভিতরের অংশ optional setting।}

% ---------------------------------------------
% SECTION 4: Environment Syntax
% ---------------------------------------------
\section{\bn{Environment Syntax}}

\bn{অনেক LaTeX structure environment দিয়ে লেখা হয়। Environment-এর শুরুতে
`\textbackslash begin\{...\}` এবং শেষে `\textbackslash end\{...\}` দিতে হয়।}

\begin{verbatim}
\begin{center}
  This text is centered.
\end{center}
\end{verbatim}

\bn{যে environment শুরু করবেন, একই environment শেষ করতেই হবে।}

% ---------------------------------------------
% SECTION 5: Special Characters
% ---------------------------------------------
\section{\bn{Special Character}}

\bn{কিছু character LaTeX-এ special meaning রাখে। এগুলো normal text হিসেবে
লিখতে চাইলে escape করতে হয়।}

\begin{center}
\begin{tabular}{ll}
\toprule
\bn{লিখতে চাইলে} & \bn{LaTeX code} \\
\midrule
\% & \verb|\%| \\
\& & \verb|\&| \\
\_ & \verb|\_| \\
\# & \verb|\#| \\
\$ & \verb|\$| \\
\{ & \verb|\{| \\
\} & \verb|\}| \\
\bottomrule
\end{tabular}
\end{center}

\bn{উদাহরণ:}

\begin{verbatim}
The model achieved 95\% accuracy.
Computer Science \& Engineering
student\_id
\end{verbatim}

% ---------------------------------------------
% SECTION 6: Heading / Section
% ---------------------------------------------
\section{\bn{Heading বা Section}}

\bn{Document structure তৈরি করতে heading ব্যবহার করা হয়।}

\begin{verbatim}
\chapter{Introduction}
\section{Background}
\subsection{Motivation}
\subsubsection{Dataset Details}
\end{verbatim}

\bn{Thesis template-এ chapter সাধারণত `main.tex` file-এ থাকে। তাই chapter file-এর
ভিতরে সাধারণত `\textbackslash section\{...\}` দিয়ে শুরু করলেই যথেষ্ট।}

% ---------------------------------------------
% SECTION 7: Paragraph and Line Break
% ---------------------------------------------
\section{\bn{Paragraph এবং Line Break}}

\bn{সাধারণ লেখা সরাসরি লিখলেই হয়।}

\begin{verbatim}
Intrusion detection is an important task in network security.
This thesis studies machine learning based detection methods.
\end{verbatim}

\bn{নতুন paragraph শুরু করতে একটি blank line দিন।}

\begin{verbatim}
This is the first paragraph.

This is the second paragraph.
\end{verbatim}

\bn{Manual line break দরকার হলে `\textbackslash\textbackslash` ব্যবহার করুন।}

\begin{verbatim}
First line\\
Second line
\end{verbatim}

% ---------------------------------------------
% SECTION 8: Text Formatting
% ---------------------------------------------
\section{\bn{Text Formatting}}

\begin{verbatim}
\textbf{bold text}
\textit{italic text}
\underline{underlined text}
\texttt{monospace code text}
\end{verbatim}

\bn{উদাহরণ:}

\begin{verbatim}
\textbf{CNN} is used for feature extraction.
\end{verbatim}

% ---------------------------------------------
% SECTION 9 (NEW): Font Size
% ---------------------------------------------
\section{\bn{Font Size}}

\bn{LaTeX-এ font size বদলাতে নিচের command গুলো ব্যবহার করা হয়।
এগুলো সবচেয়ে ছোট থেকে সবচেয়ে বড় ক্রমে সাজানো:}

\begin{verbatim}
{\tiny        এটি সবচেয়ে ছোট}
{\scriptsize  এটি খুব ছোট}
{\footnotesize এটি footnote size}
{\small       এটি একটু ছোট}
{\normalsize  এটি default size}
{\large       এটি একটু বড়}
{\Large       এটি আরও বড়}
{\LARGE       এটি বেশ বড়}
{\huge        এটি অনেক বড়}
{\Huge        এটি সবচেয়ে বড়}
\end{verbatim}

\bn{Curly brace দিয়ে আবদ্ধ করলে শুধু সেই অংশে size প্রযোজ্য হয়।}

\begin{verbatim}
Normal text and {\Large this part is larger} then back to normal.
\end{verbatim}

% ---------------------------------------------
% SECTION 10 (NEW): Text Color
% ---------------------------------------------
\section{\bn{Text Color}}

\bn{Color ব্যবহার করতে preamble-এ `xcolor` package লাগবে:}

\begin{verbatim}
\usepackage{xcolor}
\end{verbatim}

\bn{তারপর যেকোনো জায়গায়:}

\begin{verbatim}
\textcolor{red}{This text is red.}
\textcolor{blue}{This text is blue.}
\textcolor{green}{This text is green.}
\end{verbatim}

\bn{Custom color define করতে:}

\begin{verbatim}
\definecolor{myblue}{RGB}{0, 102, 204}
\textcolor{myblue}{This uses a custom color.}
\end{verbatim}

\bn{Common color নামগুলো: \texttt{red, blue, green, black, white, gray,
orange, violet, cyan, magenta, yellow, brown}.}

% ---------------------------------------------
% SECTION 11 (NEW): Spacing
% ---------------------------------------------
\section{\bn{Spacing}}

\bn{Vertical এবং horizontal space control করা যায়।}

\bn{\textbf{Vertical space} (উপর-নিচ ফাঁক):}

\begin{verbatim}
First paragraph.

\vspace{1cm}   % 1 centimeter নিচে যাবে

Second paragraph.
\end{verbatim}

\bn{\textbf{Horizontal space} (বাম-ডান ফাঁক):}

\begin{verbatim}
Word one \hspace{2cm} Word two
\end{verbatim}

\bn{\textbf{Negative space} দিয়ে কাছে আনা যায়:}

\begin{verbatim}
\vspace{-0.5cm}
\hspace{-1cm}
\end{verbatim}

\bn{Unit হিসেবে `cm`, `mm`, `pt`, `em`, `ex` ব্যবহার করা যায়।
`em` হলো current font-এর width-এর সমান।}

\bn{\textbf{Paragraph indent বন্ধ করতে:}}

\begin{verbatim}
\noindent This paragraph has no indent.
\end{verbatim}

% ---------------------------------------------
% SECTION 12 (NEW): Footnote
% ---------------------------------------------
\section{\bn{Footnote}}

\bn{Footnote হলো page-এর নিচে ছোট টীকা। এটি automatically number হয়।}

\begin{verbatim}
Machine learning\footnote{A subset of artificial intelligence.}
is used here.
\end{verbatim}

\bn{এটি compile হলে ``Machine learning'' শব্দের পরে একটি superscript number
দেখাবে এবং page-এর নিচে সেই note টি print হবে।}

% ---------------------------------------------
% SECTION 13 (NEW): Hyperlink
% ---------------------------------------------
\section{\bn{Hyperlink}}

\bn{Link ব্যবহার করতে preamble-এ `hyperref` package লাগবে:}

\begin{verbatim}
\usepackage{hyperref}
\end{verbatim}

\bn{\textbf{URL সরাসরি দেখাতে:}}

\begin{verbatim}
\url{https://www.example.com}
\end{verbatim}

\bn{\textbf{Custom text দিয়ে link বানাতে:}}

\begin{verbatim}
\href{https://www.example.com}{Visit Example Website}
\end{verbatim}

\bn{\textbf{Email link:}}

\begin{verbatim}
\href{mailto:you@example.com}{Send Email}
\end{verbatim}

\bn{Link-এর color পরিবর্তন করতে preamble-এ:}

\begin{verbatim}
\hypersetup{
  colorlinks=true,
  linkcolor=black,   % internal link color
  urlcolor=blue      % URL color
}
\end{verbatim}

% ---------------------------------------------
% SECTION 14: List
% ---------------------------------------------
\section{\bn{List}}

\bn{Bullet list:}

\begin{verbatim}
\begin{itemize}
  \item First point
  \item Second point
\end{itemize}
\end{verbatim}

\bn{Numbered list:}

\begin{verbatim}
\begin{enumerate}
  \item Collect dataset
  \item Preprocess data
  \item Train model
\end{enumerate}
\end{verbatim}

% ---------------------------------------------
% SECTION 15: Figure
% ---------------------------------------------
\section{\bn{Figure Add করা}}

\bn{Image আগে project-এর `figures/` folder-এ রাখুন। তারপর:}

\begin{verbatim}
\begin{figure}[H]
  \centering
  \includegraphics[width=0.8\textwidth]{figures/example.png}
  \caption{Example result}
  \label{fig:example-result}
\end{figure}
\end{verbatim}

\bn{Figure reference করতে:}

\begin{verbatim}
Figure~\ref{fig:example-result} shows the result.
\end{verbatim}

\bn{`\textbackslash label\{...\}` হলো reference key. একই key দিয়ে
`\textbackslash ref\{...\}` করলে figure number পাওয়া যায়।}

% ---------------------------------------------
% SECTION 16: Table
% ---------------------------------------------
\section{\bn{Table Add করা}}

\begin{verbatim}
\begin{table}[H]
  \centering
  \caption{Example table}
  \label{tab:example}
  \begin{tabular}{ll}
    \toprule
    Item & Value \\
    \midrule
    Accuracy & 95\% \\
    Loss & 0.12 \\
    \bottomrule
  \end{tabular}
\end{table}
\end{verbatim}

\bn{Column format:}

\begin{itemize}
  \item \code{\{ll\}} \bn{মানে দুইটি left aligned column}
  \item \code{\{lc\}} \bn{মানে first column left, second column center}
  \item \code{\{lcr\}} \bn{মানে left, center, right}
\end{itemize}

% ---------------------------------------------
% SECTION 17: Equation
% ---------------------------------------------
\section{\bn{Equation}}

\bn{Inline equation:}

\begin{verbatim}
Accuracy is calculated as $Accuracy = \frac{TP + TN}{TP + TN + FP + FN}$.
\end{verbatim}

\bn{Displayed equation:}

\begin{verbatim}
\begin{equation}
Accuracy = \frac{TP + TN}{TP + TN + FP + FN}
\end{equation}
\end{verbatim}

\bn{Equation label সহ:}

\begin{verbatim}
\begin{equation}
F1 = \frac{2 \times \mathrm{Precision} \times \mathrm{Recall}}
{\mathrm{Precision} + \mathrm{Recall}}
\label{eq:f1-score}
\end{equation}
\end{verbatim}

% ---------------------------------------------
% SECTION 18 (NEW): Math Symbols
% ---------------------------------------------
\section{\bn{Math Symbols}}

\bn{Preamble-এ `amsmath` package add করুন:}

\begin{verbatim}
\usepackage{amsmath}
\end{verbatim}

\bn{\textbf{Greek letters} (সব math mode `\$...\$`-এর ভিতরে):}

\begin{verbatim}
$\alpha$   $\beta$   $\gamma$   $\delta$
$\theta$   $\lambda$ $\mu$      $\sigma$
$\pi$      $\omega$  $\Sigma$   $\Delta$
\end{verbatim}

\bn{\textbf{Common math operators:}}

\begin{verbatim}
$\sum_{i=1}^{n} x_i$        % summation
$\prod_{i=1}^{n} x_i$       % product
$\int_0^\infty f(x) dx$     % integral
$\lim_{x \to \infty} f(x)$  % limit
$\sqrt{x}$                  % square root
$\sqrt[3]{x}$               % cube root
$x^{2}$                     % superscript
$x_{i}$                     % subscript
\end{verbatim}

\bn{\textbf{Comparison এবং logic:}}

\begin{verbatim}
$\leq$   % less than or equal
$\geq$   % greater than or equal
$\neq$   % not equal
$\approx$ % approximately equal
$\in$    % element of
$\notin$ % not element of
$\cup$   % union
$\cap$   % intersection
$\infty$ % infinity
\end{verbatim}

\bn{\textbf{Matrix:}}

\begin{verbatim}
\begin{equation}
A = \begin{bmatrix}
  1 & 2 \\
  3 & 4
\end{bmatrix}
\end{equation}
\end{verbatim}

% ---------------------------------------------
% SECTION 19 (NEW): Code Listing
% ---------------------------------------------
\section{\bn{Code Listing}}

\bn{Code সুন্দরভাবে দেখাতে `listings` package ব্যবহার করা হয়।
Preamble-এ add করুন:}

\begin{verbatim}
\usepackage{listings}
\usepackage{xcolor}

\lstset{
  basicstyle=\ttfamily\small,
  keywordstyle=\color{blue},
  commentstyle=\color{gray},
  stringstyle=\color{orange},
  numbers=left,
  numberstyle=\tiny,
  breaklines=true,
  frame=single
}
\end{verbatim}

\bn{Code block লিখতে:}

\begin{verbatim}
\begin{lstlisting}[language=Python, caption={My Python Code}]
def train_model(X, y):
    model = RandomForest()
    model.fit(X, y)
    return model
\end{lstlisting}
\end{verbatim}

\bn{Inline code snippet দিতে:}

\begin{verbatim}
Use \lstinline{print("Hello")} to output text.
\end{verbatim}

\bn{Supported language-এর মধ্যে আছে: \texttt{Python, C, C++, Java, R, Matlab, SQL, bash} ইত্যাদি।}

% ---------------------------------------------
% SECTION 20 (NEW): Cross-referencing
% ---------------------------------------------
\section{\bn{Cross-referencing}}

\bn{Document-এর যেকোনো জায়গায় label দিয়ে সেটি refer করা যায়।
Figure, Table, Equation, Section সবকিছুতেই এটি কাজ করে।}

\begin{verbatim}
% Label দেওয়া
\section{Introduction}
\label{sec:intro}

\begin{figure}[H]
  ...
  \label{fig:result}
\end{figure}

\begin{table}[H]
  ...
  \label{tab:data}
\end{table}

\begin{equation}
  y = mx + c
  \label{eq:linear}
\end{equation}
\end{verbatim}

\bn{Refer করার command:}

\begin{verbatim}
See Section~\ref{sec:intro} for details.
As shown in Figure~\ref{fig:result}, the model performs well.
Table~\ref{tab:data} lists the dataset statistics.
Equation~\ref{eq:linear} represents the linear model.
\end{verbatim}

\bn{`\textasciitilde` দিয়ে non-breaking space দেওয়া হয় যাতে number
পরের লাইনে চলে না যায়।}

\bn{Page number reference করতে:}

\begin{verbatim}
Refer to page~\pageref{fig:result}.
\end{verbatim}

% ---------------------------------------------
% SECTION 21 (NEW): Abstract
% ---------------------------------------------
\section{\bn{Abstract}}

\bn{Thesis বা paper-এ abstract যোগ করতে:}

\begin{verbatim}
\begin{abstract}
This thesis proposes a machine learning based approach
for network intrusion detection. Experimental results show
that the proposed method achieves 97\% accuracy on the
benchmark dataset.
\end{abstract}
\end{verbatim}

\bn{Abstract সাধারণত `\textbackslash maketitle`-এর পরেই দেওয়া হয়,
`\textbackslash tableofcontents`-এর আগে।}

% ---------------------------------------------
% SECTION 22: Citation
% ---------------------------------------------
\section{\bn{Citation}}

\bn{Reference আগে `references.bib` file-এ add করতে হয়।}

\begin{verbatim}
@article{smith2024ids,
  author  = {Smith, John},
  title   = {Intrusion Detection Using Machine Learning},
  journal = {Example Journal},
  year    = {2024}
}
\end{verbatim}

\bn{Chapter file-এ cite করতে:}

\begin{verbatim}
Machine learning has been widely used for IDS \cite{smith2024ids}.
\end{verbatim}

\bn{Multiple citation:}

\begin{verbatim}
Several studies explored IDS using deep learning \cite{smith2024ids,doe2023cnn}.
\end{verbatim}

% ---------------------------------------------
% SECTION 23 (NEW): Bibliography Compile
% ---------------------------------------------
\section{\bn{Bibliography Compile করার পদ্ধতি}}

\bn{Citation কাজ করাতে শুধু `\textbackslash cite\{\}` লিখলেই হবে না।
`references.bib` file compile করতে হবে। সাধারণ compile sequence:}

\begin{verbatim}
lualatex main.tex
bibtex main
lualatex main.tex
lualatex main.tex
\end{verbatim}

\bn{অথবা `make` ব্যবহার করলে যদি Makefile-এ এই sequence থাকে, তাহলে
একটি command-এই হয়ে যাবে।}

\bn{Document-এর শেষে bibliography print করতে:}

\begin{verbatim}
\bibliographystyle{plain}      % citation style
\bibliography{references}      % references.bib file-এর নাম
\end{verbatim}

\bn{Common bibliography style:}

\begin{itemize}
  \item \code{plain} \bn{--- author alphabetical order, numbered}
  \item \code{ieeetr} \bn{--- IEEE format, citation order অনুযায়ী}
  \item \code{apalike} \bn{--- APA format}
\end{itemize}

% ---------------------------------------------
% SECTION 24 (NEW): Input and Include
% ---------------------------------------------
\section{\bn{\textbackslash input এবং \textbackslash include}}

\bn{বড় document (thesis, book) লেখার সময় সব কিছু একটি file-এ না লিখে
আলাদা আলাদা file-এ ভাগ করে রাখা ভালো।}

\bn{\textbf{\textbackslash input\{...\}} --- অন্য file-এর content সেখানে বসিয়ে দেয়:}

\begin{verbatim}
% main.tex-এ:
\section{Background}

Introduce the research area and explain why the problem matters.

\section{Problem Statement}

State the specific problem addressed by this thesis.

\section{Objectives}

List the main objectives of the research.

\section{Thesis Organization}

Briefly describe the structure of the remaining chapters.

Describe the methods, tools, datasets, experimental setup, or system design used
in the thesis.

Present and discuss the results of the work.

\begin{table}[H]
  \centering
  \caption{Example table}
  \label{tab:example}
  \begin{tabular}{ll}
    \toprule
    Item & Value \\
    \midrule
    Example & 1.00 \\
    \bottomrule
  \end{tabular}
\end{table}

\end{verbatim}

\bn{\textbf{\textbackslash include\{...\}} --- আলাদা page-এ শুরু করে (chapter-এর জন্য ভালো):}

\begin{verbatim}
\section{Background}

Introduce the research area and explain why the problem matters.

\section{Problem Statement}

State the specific problem addressed by this thesis.

\section{Objectives}

List the main objectives of the research.

\section{Thesis Organization}

Briefly describe the structure of the remaining chapters.

Describe the methods, tools, datasets, experimental setup, or system design used
in the thesis.

\end{verbatim}

\bn{পার্থক্য: `\textbackslash input` continuous text বসায়, কোনো page break দেয় না।
`\textbackslash include` নতুন page-এ শুরু করে এবং selective compile করা যায়।}

\bn{Selective compile (শুধু কিছু chapter দেখতে):}

\begin{verbatim}
\includeonly{chapters/introduction, chapters/results}
\end{verbatim}

% ---------------------------------------------
% SECTION 25: Comment
% ---------------------------------------------
\section{\bn{Comment}}

\bn{Percent sign `\%` দিয়ে comment লেখা হয়। Comment PDF-এ দেখাবে না।}

\begin{verbatim}
% This line will not appear in the PDF.
This line will appear in the PDF.
\end{verbatim}

% ---------------------------------------------
% SECTION 26: Bangla writing rules
% ---------------------------------------------
\section{\bn{Bangla লেখার নিয়ম}}

\bn{এই file compile করতে LuaLaTeX ব্যবহার করুন। বাংলা লিখতে এই file-এ
`\textbackslash bn\{...\}` command ব্যবহার করা হয়েছে।}

\begin{verbatim}
\bn{এখানে বাংলা লিখুন}
\end{verbatim}

\bn{বাংলা + English একসাথে লেখা যায়:}

\begin{verbatim}
\bn{এটি বাংলা লেখা।} This is English text.
\end{verbatim}

% ---------------------------------------------
% SECTION 27: Compile Command
% ---------------------------------------------
\section{\bn{Compile Command}}

\bn{এই file compile করতে `la-tex basic` folder-এ গিয়ে চালান:}

\begin{verbatim}
make
\end{verbatim}

\bn{Bangla font ঠিকমতো না এলে `Noto Serif Bengali` এবং `Noto Sans Bengali`
install আছে কিনা check করুন।}

\end{document}
